\documentclass[10pt,a4paper]{amsart}
\usepackage[margin=19mm]{geometry}
\usepackage{amsmath,amssymb,mathtools}
\usepackage[scr=rsfs,cal=euler]{mathalfa}
\usepackage{xfrac}
\usepackage[unicode,hidelinks,bookmarks=false]{hyperref}
\newcommand{\ie}{i.e.\ }
\newcommand{\cf}{cf.\ }
\newcommand{\resp}{respectively\ }
\newcommand{\Subsection}{\subsection}
\providecommand{\nobreakdash}{\nobreak\hbox{-}}
\setlength{\emergencystretch}{2em}
\multlinegap0pt
\usepackage{bm}
\usepackage{bbm}
\usepackage{dsfont}
\usepackage{xspace}
\usepackage{cancel}
\usepackage{mathtools}
\usepackage{cleveref}
\usepackage{amsthm}
\makeatletter
\@ifundefined{definition}{\newtheorem{definition}{Definition}}{}
\@ifundefined{lemma}{\newtheorem{lemma}{Lemma}}{}
\@ifundefined{proposition}{\newtheorem{proposition}{Proposition}}{}
\@ifundefined{remark}{\newtheorem{remark}{Remark}}{}
\makeatother
\newcommand*{\defined}{\coloneqq}
\newcommand*{\wealth}{\mathcal{K}}
\newcommand{\X}{\mathcal{X}}
\newcommand{\dG}{d_{\mc{G}}} 
\newcommand{\dkl}{d_{\mathrm{KL}}}
\newcommand{\iid}{i.i.d.\xspace}
\newcommand{\lb}{\left [}
\newcommand{\lp}{\left (} %
\newcommand{\mc}[1]{\mathcal{#1}}
\newcommand{\prediction}{\mc{A}_{\text{pred}}}
\newcommand{\rb}{\right ]}
\newcommand{\rp}{\right )} %
\title[Nonparametric Two-Sample Testing by Betting]{Nonparametric Two-Sample Testing by Betting}
\author{Shubhanshu Shekhar, Aaditya Ramdas}
\date{Source: arXiv:2112.09162v6 (CC BY 4.0)}
\begin{document}
\maketitle

\noindent\emph{Source and licence.} Nonparametric Two-Sample Testing by Betting. Version arXiv:2112.09162v6 (CC BY 4.0), distributed under the Creative Commons Attribution 4.0 International licence, as recorded in the arXiv licence metadata for this exact version and shown on the abstract page https://arxiv.org/abs/2112.09162v6. Full rights notice: licenses/arxiv-2112.09162-CC-BY-4.0.txt.\par\smallskip

\noindent\textit{Editorial scope.} Counted unit: the Remark whose source label is \texttt{None} in the article (source0.tex lines 17346--17348, source label \texttt{None}), delivered together with the article's own complete original proof of it (source0.tex lines 17349--17381, one \texttt{proof} environment of 33 lines). No mathematics is written, repaired or re-derived: statement and proof are the article's own text line for line. Required context retained uncounted: def:ONS (lines 447--462); eq:oga-1 (lines 633--636); def:bounded-means-test (lines 640--649); prop:bounded-mean (lines 658--665); prop:kernel-mmd (lines 730--738); lemma:lower-bound-general (lines 17340--17345). Every definition, display and label of those lines is retained unchanged. Omitted from this package and not redistributed: 1--446, 463--632, 637--639, 650--657, 666--729, 739--17339, 17382--17786. Nothing inside a retained excerpt is omitted or altered: every retained body is delivered exactly as the article's lines are written, so no omission receipt of any kind accompanies this package. Citations: the retained excerpts carry no citation command at all, so no bibliography entry of the article is transcribed and no reference is redistributed. Reference adaptation: none was needed and none was made; every reference inside the delivered unit resolves inside it, so no unresolved reference or citation is produced. Printed slips kept literal: no printed word, symbol, hypothesis or number of the counted statement or of its proof was altered. Layout and class: the article's own class is \texttt{article} with packages including geometry, layout, tikz, xcolor, wrapfig, soul, url, amsmath, amsthm, amssymb, bm, bbm, mathtools, comment; the wrapper is the lane's pinned \texttt{amsart} editorial wrapper at 10pt on A4, which restates the theorem-like environments the retained text needs and adds the article's own macro definitions that the retained text needs (\texttt{\textbackslash X}, \texttt{\textbackslash dG}, \texttt{\textbackslash defined}, \texttt{\textbackslash dkl}, \texttt{\textbackslash iid}, \texttt{\textbackslash lb}, \texttt{\textbackslash lp}, \texttt{\textbackslash mc}, \texttt{\textbackslash prediction}, \texttt{\textbackslash rb}, \texttt{\textbackslash rp}, \texttt{\textbackslash wealth}), with extra wrapper packages (bm, bbm, dsfont, xspace, cancel, mathtools, cleveref). The delivered numbering is the unit's own. Assets: no retained span contains a figure environment, a graphics-inclusion command, a PDF-page inclusion, a table or a TikZ picture; no image file is copied into this package, the package declares no asset dependency and no image file is redistributed with it. Licence: version arXiv:2112.09162v6 of this article is distributed under the Creative Commons Attribution 4.0 International licence, as recorded in the arXiv licence metadata for this exact version and shown on the abstract page https://arxiv.org/abs/2112.09162v6. A full rights notice is delivered at licenses/arxiv-2112.09162-CC-BY-4.0.txt. Characters: every retained excerpt is pure ASCII, so the delivered engine, XeLaTeX with its default Latin Modern text font, reports no missing character.
    \begin{definition}[ONS betting strategy]
        \label{def:ONS}
        Let $\{v_t \in [-1,1]: t \geq 1\}$ denote a sequence of outcomes. 
        In the two-sample testing case, we have $v_t = g_t(X_t) - g_t(Y_t)$, where $\{g_t: t \geq 1\}$ are chosen using any prediction strategy $\prediction$.
        Initialize $\lambda_1 = 0$, and $a_0=1$. Then, for $t=1, 2, \ldots:$
        \begin{itemize}
            \item Observe $v_t \in [-1,1]$. 
            \item Set $z_t = {v_t / (1 + v_t \lambda_t)}$. 
            \item Update $a_t = a_{t-1} + z_t^2$. 
            \item Update $\lambda_{t+1}$ as follows: 
            \begin{align}
                \label{def:lambda-ons}
                \lambda_{t+1} = \min \left\{\frac{1}{2}, \max \left\{ - \frac{1}{2}, {\lambda_t + \frac{2}{2-\log 3} \frac{z_t}{a_t} }\right\} \right\}. 
            \end{align}
        \end{itemize}
    \end{definition}

        \begin{align}
            \label{eq:oga-1}
            u_{t+1} =  \Pi_{\mc{U}} \lp u_{t} + \frac{1}{\sqrt{M_t}}\lp X_{t} - Y_{t} \rp \rp,
        \end{align}

        \begin{definition}[Sequential Test for Bounded Means Testing]
        \label{def:bounded-means-test}
            Set $\wealth_0 = 1$, $\lambda_1=0$, and $u_1=\boldsymbol{0} \in \mathbb{R}^m$. For $t=1, 2, \ldots$:
            \begin{itemize}
                \item Observe $X_t, Y_t$. 
                \item Update the wealth: $\wealth_{t} = \wealth_{t-1} \times \lp 1 + \lambda_t \langle u_t, X_t - Y_t \rangle \rp$. Reject the null if $\wealth_t \geq 1/\alpha$. 
                \item Obtain $u_{t+1}$ from the OGA prediction strategy~\eqref{eq:oga-1}. 
                \item Obtain $\lambda_{t+1}$ from the ONS betting strategy of~\Cref{def:ONS}. 
            \end{itemize}
        \end{definition}

    \begin{proposition}
    \label{prop:bounded-mean}
        Introduce the terms $\Delta = \dG(P_X, P_Y) =  (1/2)\|\mathbb{E}_{P_X}[X] - \mathbb{E}_{P_Y}[Y] \|_2$,  and  $\sigma^2 = \sup_{u:\|u\|\leq 1/2} \mathbb{V}\lp \langle u, X-Y\rangle \rp = \mc{O}\lp \mathbb{E}[\|X-Y\|^2]\rp$. Then, for the sequential test described in~\Cref{def:bounded-means-test}, we have the following: 
        \begin{align}
            &\text{Under } H_0: \; \mathbb{P}(\tau < \infty) \leq \alpha. \\
            &\text{Under } H_1: \; \mathbb{P} \lp \tau <\infty \rp = 1, \quad \text{and} \quad \mathbb{E}[\tau] = \mc{O} \lp \frac{ \sigma^2 \log(1/\Delta \alpha) }{\Delta^2}  {+ \frac{\log(1/\alpha\Delta)}{\Delta}}\rp. \label{eq:tau-bounded-mean-1}
        \end{align}
    \end{proposition}

            \begin{proposition}
            \label{prop:kernel-mmd}
                For the sequential test described in~\Cref{def:bounded-means-test}, we have the following: 
                \begin{align}
                    \text{Under } H_0: \; &\mathbb{P}(\tau < \infty) \leq \alpha. \\
                    \text{Under } H_1: \; &\mathbb{P} \lp \tau <\infty \rp = 1,  \quad \mathbb{E}[\tau] = \mc{O}\lp  \frac{\sigma^2  \log(1/\alpha \Delta)}{\Delta^2}  +{\frac{\log(1/\alpha \Delta)}{\Delta}}\rp,  \\ 
                    \text{and} \quad &\lim_{n \to \infty} -\frac{1}{2n} \log   \mathbb{P} \lp \tau > n \rp   \geq \inf_{P' \in \mc{P}(\X)} \; \frac{1}{2} \lp \dkl\lp P', P_X \rp + \dkl \lp P', P_Y \rp \rp. 
                \end{align}
            \end{proposition}

        \begin{lemma}
            \label{lemma:lower-bound-general} Given a stream of observations $Z_1, Z_2, \ldots \stackrel{\iid}{\sim} P_Z$, consider the problem of testing the null $H_0: P_Z = Q_0$ against the alternative $H_1: P_Z = Q_1$, for $Q_0 \neq Q_1$. Let $\tau'$ denote any level-$\alpha$, power-one, sequential test for this problem. Then, we have 
            \begin{align}
                \mathbb{E}_{H_1}[\tau'] \geq \frac{\log(1/\alpha)}{\dkl(Q_1, Q_0)}. 
            \end{align}
        \end{lemma}

        \begin{remark}
            To apply this result to the bounded-mean testing and two-sample testing problems, we will construct distributions $P_X$ and $P_Y$, and set $Q_0 = P_X \times P_X$, and $Q_1 = P_X \times P_Y$. Then, for any level-$\alpha$ sequential test for these problems,~\Cref{lemma:lower-bound-general} immediately implies that $\mathbb{E}_{H_1}[\tau'] \geq \log(1/\alpha)/\dkl(P_Y, P_X)$. With an appropriate choice of the distributions $P_X$ and $P_Y$, we get lower bound expressions matching the upper bounds in~\Cref{prop:bounded-mean} and~\Cref{prop:kernel-mmd}.
        \end{remark}

        \begin{proof}
         The proof of this statement uses some standard properties of the KL-divergence that are often used in proving lower bounds for multi-armed bandit algorithms, as we detail below.  

        Since $\tau'$ is a power-one test with finite expectation under $H_1$, we can use Wald's Lemma to get 
        \begin{align}
            \dkl\lp Q_1^{\otimes \tau'}, Q_0^{\otimes \tau'} \rp &= \mathbb{E}_{H_1} \lb \log \lp \frac{dQ_1^{\otimes \tau'}}{dQ_0^{\otimes \tau'}} \rp \rb
             = \mathbb{E}_{H_1} \lb \sum_{t=1}^{\tau'} \log \lp \frac{dQ_1}{dQ_0} \rp \rb \\
             & = \mathbb{E}_{H_1}[\tau'] \dkl(Q_1, Q_0). \label{eq:lower-general-1}
        \end{align}
        Note that the last equality used the assumptions that $\tau'$ has finite expectation under the alternative, and that $0<\dkl(Q_1, Q_0)<\infty$. 

        Next, we consider the event $E = \{\tau' < \infty\}$. By definition of a sequential level-$\alpha$ test of power one, we have 
        \begin{align}
            q_1 \defined \mathbb{P}_{H_1}\lp E \rp = 1, \quad \text{and} \quad 
            q_0 \defined \mathbb{P}_{H_0} \lp E \rp \leq \alpha. 
        \end{align}

        Next, on evaluating KL divergence between two Bernoulli distributions with parameters $q_1$ and $q_0$, we get
        \begin{align}
            \dkl(q_1, q_0) &= q_1 \log \lp \frac{q_1}{q_0} \rp + (1-q_1) \log \lp \frac{1-q_1}{1-q_0} \rp \\
            & = \log \lp \frac{1}{q_0} \rp \geq \log \lp \frac{1}{\alpha} \rp. \label{eq:lower-general-2}
        \end{align}

        The final step is to relate~\eqref{eq:lower-general-1} and~\eqref{eq:lower-general-2} by using the data-processing inequality, as follows 
        \begin{align}
            \dkl(Q_1^{\otimes \tau'}, Q_0^{\otimes \tau'}) = \mathbb{E}_{H_1}[\tau'] \dkl(Q_1, Q_0) \geq \dkl(q_1, q_0),  
        \end{align}
        which implies the required statement
        \begin{align}
            \mathbb{E}_{H_1}[\tau'] \geq \frac{\log(1/\alpha)}{\dkl(Q_1, Q_0)}.
        \end{align}
        This completes the proof of~\Cref{lemma:lower-bound-general}. 
        \end{proof}

\end{document}
